%==============================================================================
% Chapitre 7 - Température
%==============================================================================

\section{Introduction}

La température est l'une des grandeurs physiques les plus importantes en
balistique.

Elle influence :

\begin{itemize}
\item la densité de l'air ;
\item la vitesse du son ;
\item le comportement des poudres propulsives ;
\item les performances des systèmes de mesure.
\end{itemize}

Ses effets sont parfois faibles à courte distance mais deviennent
significatifs dès que la portée augmente.

\begin{aretenir}

Deux tirs réalisés avec la même arme et la même munition peuvent produire des
résultats différents si la température change.

\end{aretenir}

\section{Qu'est-ce que la température ?}

La température caractérise l'état thermique d'un système.

À l'échelle microscopique, elle est liée à l'agitation des molécules qui
composent la matière.

Plus la température est élevée :

\begin{itemize}
\item plus les molécules sont agitées ;
\item plus elles occupent d'espace ;
\item plus certaines propriétés physiques évoluent.
\end{itemize}

Dans l'atmosphère, cette agitation influence directement le comportement de
l'air.

\section{Unités de température}

Dans la vie courante, la température est généralement exprimée en degrés
Celsius (°C).

Dans les calculs physiques, on utilise souvent la température absolue,
exprimée en kelvins (K).

La conversion est donnée par :

\begin{equation}
T_K = T_{^\circ C} + 273,15
\end{equation}

où :

\begin{itemize}
\item $T_K$ est la température en kelvins ;
\item $T_{^\circ C}$ est la température en degrés Celsius.
\end{itemize}

\section{Température et densité de l'air}

Lorsque l'air se réchauffe, il se dilate.

À pression égale :

\begin{itemize}
\item son volume augmente ;
\item sa densité diminue.
\end{itemize}

Inversement, un air froid est généralement plus dense.

Cette variation de densité influence directement la traînée exercée sur les
projectiles.

\begin{exemple}

Un tir réalisé à 35°C rencontrera généralement moins de résistance
aérodynamique qu'un tir identique réalisé à 0°C.

\end{exemple}

\section{Température et vitesse du son}

La vitesse du son dépend de la température de l'air.

Une approximation couramment utilisée est :

\begin{rigueur}

\begin{equation}
c \approx 331 + 0,6,T
\label{eq:vitesse_son_temperature}
\end{equation}

où :

\begin{itemize}
\item $c$ est la vitesse du son en m/s ;
\item $T$ est la température en °C.
\end{itemize}

\end{rigueur}

\begin{exemple}

À 0°C :

\begin{equation}
c \approx 331~\mathrm{m/s}
\end{equation}

À 20°C :

\begin{equation}
c \approx 343~\mathrm{m/s}
\end{equation}

\end{exemple}

Cette variation influence directement le nombre de Mach, qui sera étudié dans
un chapitre ultérieur.

\section{Température et munitions}

La température n'influence pas uniquement l'atmosphère.

Elle agit également sur les munitions.

Les poudres propulsives peuvent présenter des comportements différents selon
la température.

De manière générale :

\begin{itemize}
\item une température élevée tend à augmenter l'énergie disponible ;
\item une température faible tend à la diminuer.
\end{itemize}

L'importance de cet effet dépend fortement :

\begin{itemize}
\item de la formulation de la poudre ;
\item du type de munition ;
\item des écarts de température considérés.
\end{itemize}

\section{Température des armes}

La température de l'arme elle-même peut également évoluer.

Lors de tirs répétés :

\begin{itemize}
\item le canon s'échauffe ;
\item certaines dimensions évoluent légèrement ;
\item les conditions de fonctionnement peuvent être modifiées.
\end{itemize}

Ces phénomènes sont particulièrement importants pour :

\begin{itemize}
\item les armes automatiques ;
\item les essais d'endurance ;
\item les tirs de précision.
\end{itemize}

\section{Influence sur les essais}

Lors d'une campagne d'essais, la température doit être relevée
systématiquement.

Cette mesure permet :

\begin{itemize}
\item d'interpréter certains écarts ;
\item de comparer des résultats obtenus à des dates différentes ;
\item d'améliorer la corrélation entre essais et simulation.
\end{itemize}

\begin{simulation}

Les moteurs balistiques avancés prennent en compte la température afin
d'ajuster :

\begin{itemize}
\item la densité de l'air ;
\item la vitesse du son ;
\item certains paramètres aérodynamiques.
\end{itemize}

L'influence peut être faible à courte portée mais devient notable à longue
distance.

\end{simulation}

\section{Application à la balistique}

La température influence simultanément :

\begin{itemize}
\item le milieu traversé par le projectile ;
\item les performances de la munition ;
\item les conditions d'essais ;
\item certains paramètres de simulation.
\end{itemize}

Elle constitue donc un paramètre environnemental essentiel.

\section{À retenir}

\begin{aretenir}

\begin{itemize}
\item La température influence la densité de l'air.
\item Elle modifie la vitesse du son.
\item Elle peut affecter les performances des munitions.
\item Elle doit être relevée lors des essais balistiques.
\item Son influence augmente généralement avec la distance de tir.
\end{itemize}

\end{aretenir}

